% Chapter number style: boxed (default) or plain
\documentclass[boxednum]{chapterstart}
% \documentclass[plainnum]{chapterstart}

\begin{document}

\chapteropener{1}{Functions and Limits}{%
Calculus is fundamentally different from the mathematics that you have
studied previously. Calculus is less static and more dynamic. It is
concerned with change and motion; it deals with quantities that approach
other quantities. So in this first chapter we begin our study of calculus
by investigating how the values of functions change and approach limits.
}

\sectionheader{1.1}{Functions and Their Representations}

Functions arise whenever one quantity depends on another. Consider the
following four situations.

\textbf{A.} The area of a circle depends on the radius of the circle. The
rule that connects $r$ and $A$ is given by the equation $A = \pi r^2$. With
each positive number $r$ there is associated one value of $A$, and we say
that $A$ is a function of $r$.
\sidenote{Input: radius $r$.\\Output: area $A$.}

\textbf{B.} The human population of the world depends on the time $t$. The
table gives estimates of the world population $P(t)$ at time $t$ for certain
years. For instance, $P(1950) = 2{,}560{,}000{,}000$. But for each value of
the time $t$ there is a corresponding value of $P$, and we say that $P$ is a
function of $t$.
\sidenote{Full data in Table~\ref{tab:population}.}

% Table in the left margin, beside the text (sample style).
% For a table that overhangs into the right margin instead, use
% \begin{widetable}[place=right, extend=1.0in] ... \end{widetable}
\begin{widetable}[place=margin]
  \caption{World population estimates, 1900--2000.}
  \label{tab:population}
  \begin{tabular*}{\linewidth}{@{\extracolsep{\fill}}lr}
    \toprule
    Year & \shortstack{Population\\(millions)} \\
    \midrule
    1900 & 1650 \\
    1910 & 1750 \\
    1920 & 1860 \\
    1930 & 2070 \\
    1940 & 2300 \\
    1950 & 2560 \\
    1960 & 3040 \\
    1970 & 3710 \\
    1980 & 4450 \\
    1990 & 5280 \\
    2000 & 6080 \\
    \bottomrule
  \end{tabular*}
\end{widetable}

\textbf{C.} The cost of mailing a first-class letter depends on the weight
$w$ of the letter. Although there is no simple formula that connects $w$ and
$C$, the post office has a rule for determining $C$ when $w$ is known.
\sidenote{A rule, not a formula.}

\textbf{D.} The vertical acceleration of the ground as measured by a
seismograph during an earthquake is a function of the elapsed time $t$.
Figure~\ref{fig:seismo} shows a graph generated by seismic activity during
the Northridge earthquake that shook Los Angeles in 1994. For a given value
of $t$ the graph provides a corresponding value of $a$.

\chapterfigure[fig:seismo]{seismograph.pdf}{%
Vertical ground acceleration during the Northridge earthquake
(dummy data for illustration).}

Each of these examples describes a rule whereby, given a number ($r$, $t$,
$w$, or $t$), another number ($A$, $P$, $C$, or $a$) is assigned. In each
case we say that the second number is a function of the first number.
\sidenote{Each case: one input, one output.}

\end{document}
